\chapter{Publishing Cartridges on the Cartridge Store}
\label{ch:store}

A game that only its author can run is an exercise. A game that a classmate
can find, download onto their own board, and beat your high score on is a
publication. The \textbf{PRG32 Cartridge Store} is the service that turns one
into the other: a catalogue of \fname{.prg32} cartridges with an editorial
review step, plus the score, metrics, and multiplayer services that games use
at run time \citep{prg32store}. This chapter follows a cartridge from a source
file to a catalogue page, and it tells the story four times, once for each kind
of person who meets the Store: the visitor without an account, the registered
user who publishes, the editor who reviews, and the administrator who runs the
service.

\section{The Store in one picture}

\begin{center}
\begin{tikzpicture}[font=\sffamily\scriptsize,node distance=5mm,
  s/.style={rounded corners=2pt,draw=prgsteel,fill=prgpanel,minimum height=9mm,minimum width=25mm,align=center},
  q/.style={rounded corners=2pt,draw=prgamber,fill=prgamber!10,minimum height=9mm,minimum width=25mm,align=center},
  g/.style={rounded corners=2pt,draw=prggreen,fill=prggreen!10,minimum height=9mm,minimum width=25mm,align=center},
  ar/.style={-{Stealth[length=2mm]},draw=prgink!70}]
  \node[s] (dev) {registered user\\builds a bundle};
  \node[q,right=9mm of dev] (pend) {pending\\submission};
  \node[s,below=6mm of pend] (ed) {editor\\verifies or rejects};
  \node[g,right=9mm of pend] (cat) {public\\catalogue};
  \node[s,right=9mm of cat,yshift=8mm] (brd) {\prg board,\\QEMU, PRG32-QT};
  \node[s,right=9mm of cat,yshift=-8mm] (web) {anyone with\\a browser};
  \draw[ar] (dev) -- node[above]{upload} (pend);
  \draw[ar] (ed) -- (pend);
  \draw[ar] (pend) -- node[above]{verify} (cat);
  \draw[ar] (cat) -- (brd);
  \draw[ar] (cat) -- (web);
\end{tikzpicture}
\captionof{figure}{The life of a submission. Nothing reaches the public
catalogue without an editor's decision \citep{prg32storeapi}.}
\label{fig:storeflow}
\end{center}

\begin{prgconcept}
Publishing is \emph{moderated}. An upload creates a pending submission; only
when a member of the \code{editors} group verifies it does the cartridge appear
in the catalogue. The Store never executes uploaded code: it validates the
package, stores it, and serves it.
\end{prgconcept}

\begin{center}
\begin{tabularx}{\linewidth}{@{}lXl@{}}
\toprule
Who & May & Needs \\
\midrule
Visitor & browse, search, read game pages and colophons, download cartridges,
  view scoreboards and statistics & nothing \\
Registered user & everything above, plus upload bundles for review, create
  publication tokens, submit scores, own metrics runs & verified e-mail \\
Editor & everything above, plus review the queue, edit descriptive metadata,
  verify or reject, administer published cartridges & \code{editors} group \\
Administrator & everything above, plus settings, users, groups, roles,
  authentication providers, backup and restore & admin role \\
\bottomrule
\end{tabularx}
\captionof{table}{The four points of view of this chapter
\citep{prg32store,prg32storeapi}.}
\label{tab:storeroles}
\end{center}

\section{What exactly gets published}

Three nested things are involved, and beginners confuse them.

A \textbf{cartridge} is one \fname{.prg32} file: the \code{PRG2} header, your
linked code and data, an optional audio block, and an optional metadata
trailer \citep{prg32cartfmt}. One cartridge file contains one linked image, so
the Store keeps one file per \emph{architecture}: \code{esp32c6} for the
physical board and \code{qemu} for the emulator.

A \textbf{bundle} is a ZIP archive holding a \fname{manifest.json}, an icon,
and one or more cartridge files. The bundle is the unit of submission.

A \textbf{game} in the catalogue is identified by the manifest's \code{id};
each \code{version} of that game holds its architecture variants.

\lstinputlisting[style=json,
  caption={A bundle manifest (\fname{examples/store\_publishing/manifest.json}).
  Required fields are \code{abi}, \code{id}, \code{title}, \code{version},
  \code{assets.icon}, and a non-empty \code{architectures} list.}
]{examples/store_publishing/manifest.json}

\begin{prgwarn}
The \code{id} and the \code{version} are forever. Editors can correct a title
or a tag during review, but the Store preserves the identifier, the version,
and the authorship of the package as submitted. Choose a reverse-domain
identifier you control (\code{it.myschool.class3b.lemons}), and raise the
version for every new submission.
\end{prgwarn}

A cartridge may also carry a \textbf{colophon}: the credits and controls screen
shown after a cartridge is activated and before play begins. It is the game's
title page, and writing one is a small lesson in attribution.

\lstinputlisting[style=json,
  caption={A colophon (\fname{examples/store\_publishing/colophon.json}).
  \code{abi}, \code{title}, \code{version}, and \code{developer.name} are
  required.}
]{examples/store_publishing/colophon.json}

\section{The visitor: no account needed}

Everything about \emph{getting} games is open. A visitor to the Store's web
page sees the catalogue, can search by title, summary, author, or tag, and can
open a game's page to read its description, authors, versions, download count,
colophon, and leaderboard. The \kbtn{Scores} and \kbtn{Stats} pages are public
too.

Most visitors, however, are not people at a browser but \emph{devices}. A \prg
board whose Wi\nobreakdash-Fi is configured finds the Store (by a saved address
or by mDNS discovery on the classroom network), lists the catalogue from its
Setup menu, and downloads a cartridge straight into one of its four slots. The
PRG32-QT application of Chapter~\ref{ch:qt} does the same from its Store
screen. Both ask for the variant that matches their architecture.

\lstinputlisting[style=bash,
  caption={The visitor's view from a terminal: discover, list, search, and
  download, all without credentials.}
]{scripts/store/client_download.sh}

\begin{center}
\begin{tabularx}{\linewidth}{@{}lX@{}}
\toprule
Endpoint & Returns \\
\midrule
\code{GET /.well-known/prg32-store.json} & the discovery document: where each
  service lives \\
\code{GET /api/games} & the catalogue; accepts \code{q}, \code{page},
  \code{per\_page} \\
\code{GET /api/games/<id>} & one game record with its versions \\
\code{GET /api/games/<id>/icon}, \code{/screenshot}, \code{/colophon} & the
  game's artwork and credits \\
\code{GET /api/games/<id>/download} & a cartridge; accepts \code{version} and
  \code{architecture} \\
\code{GET /api/scores} & scores; accepts \code{game}, \code{player},
  \code{limit} \\
\code{GET /api/stats/downloads} & download time series and top games \\
\bottomrule
\end{tabularx}
\captionof{table}{The public, read-only interface \citep{prg32storeapi}.}
\end{center}

\begin{prgaside}
The firmware and the host tools check a downloaded cartridge before running
it: the ABI major version, the ABI hash, and the features it requires must be
ones the runtime provides. An incompatible cartridge is refused with a message
rather than executed. Openness about downloading is safe because the consumer,
not the catalogue, has the last word.
\end{prgaside}

\section{The registered user: publishing your game}

\subsection{An account and a token}

To register, enter an e\nobreakdash-mail address on the Store's
\kbtn{Register} page; the Store sends a verification link, and following it
lets you set a password. Your username is that verified address. A school
Store may instead be connected to the institution's single sign-on, in which
case you simply log in.

Publishing from a command line, or from the Construction Kit, needs a
\textbf{publication token}. Open \kbtn{Tokens}, type a label that says where
the token will live (``lab laptop''), and create it. The token is displayed
once: copy it then. You can rename or revoke any of your tokens later, and you
can see when each was last used.

\begin{prgwarn}
A token \emph{is} your account, for anything a script can do. Keep it in an
environment variable, never in a source file, and never commit it to Git
(Appendix~\ref{app:github}). If one leaks, revoke it on the \kbtn{Tokens} page
and create another.
\end{prgwarn}

\subsection{Build one cartridge per architecture}

From a \prg checkout with ESP-IDF active, build the portable cartridge twice,
once for each architecture label. The source here is the C Lemon Catcher of
Appendix~\ref{app:worked}.

\lstinputlisting[style=bash,
  caption={Building the two architecture variants and inspecting one.}
]{scripts/store/build_variants.sh}

Before you submit, \emph{play} both variants: the \code{qemu} one in QEMU or
PRG32-QT, the \code{esp32c6} one on a board if you have it. An editor will.

\subsection{Describe it}

Put \fname{manifest.json} and a square PNG icon beside the two cartridge files.
Optionally embed the same metadata, the icon, and the colophon inside each
cartridge as a metadata trailer, so that the file describes itself even when it
travels outside a bundle.

\lstinputlisting[style=bash,
  caption={Optional: attaching a metadata trailer and reading it back.}
]{scripts/store/attach_metadata.sh}

\subsection{Pack and submit}

\lstinputlisting[style=bash,
  caption={Packing the bundle and submitting it for review.}
]{scripts/store/pack_and_publish.sh}

The same bundle can be submitted from a browser: log in, open \kbtn{Publish},
choose the ZIP file, and press \kbtn{Upload for Review}. Either way the Store
answers with a \emph{pending} submission and its number. Your game is not yet
public. When an editor verifies it, it appears in the catalogue; if it is
rejected, fix what the editor asked for, raise the version if required, and
submit again.

\begin{center}
\begin{tabularx}{\linewidth}{@{}lX@{}}
\toprule
Before submitting, check that\dots & Why \\
\midrule
both variants build and run & the editor will try them \\
\code{cartridge summary} shows the portable import model & firmware-specific
  cartridges are not accepted by current tools \\
the \code{id} is yours and the \code{version} is new & both are permanent \\
the title and summary say what the game is & they are the catalogue entry \\
authors and licence are true & attribution is part of the record \\
the colophon lists the controls & players read it before they play \\
no personal data is in any text & publication is public \\
all art and sound are your own or properly licensed & you are the publisher \\
\bottomrule
\end{tabularx}
\captionof{table}{A submission checklist.}
\label{tab:storechecklist}
\end{center}

\subsection{After publication: scores and runs}

A published game can use the Store at run time. Scores posted with your session
or token appear on the game's leaderboard and on the public scoreboard, and a
cartridge can submit them itself with
\code{prg32\_score\_submit\_current\_player}. Performance runs
(Appendix~\ref{app:perf}) that are posted with your token are linked to your
account and listed on your own runs page, which only you and the administrators
can open.

\section{The editor: reviewing submissions}

Editors are ordinary users whom an administrator has added to the
\code{editors} group. They see one more item in the navigation bar,
\kbtn{Review}, which opens the queue of pending submissions.

Opening a submission shows who submitted it, its identifier, version, authors,
and architecture files, and an \emph{Editable Metadata} form. An editor may
change exactly seven descriptive fields: title, summary, description, tags,
licence, homepage, and repository. The identifier, the version, and the authors
are shown but cannot be changed; an attempt to change them through the API is
refused \citep{prg32storeapi}. Two buttons end the review: \kbtn{Verify and
Publish}, and \kbtn{Reject}.

\lstinputlisting[style=bash,
  caption={The editor's review, scripted. The token belongs to a user in the
  \code{editors} group.}
]{scripts/store/editor_review.sh}

\begin{prgconcept}
The editor may improve how a game is \emph{described}, but may not alter what
it \emph{is} or who made it. This is the same separation a journal observes
between copy-editing a paper and changing its authors.
\end{prgconcept}

What should an editor actually check? The Store validates form: that the
package is a ZIP, that the manifest and colophon are well-formed, that the icon
is a real image, that identifiers are path-safe. Judgement is left to people.

\begin{center}
\begin{tabularx}{\linewidth}{@{}lX@{}}
\toprule
Question & How to answer it \\
\midrule
Does it run? & obtain each variant of the submission and play it in
  PRG32-QT or QEMU \\
Is it what it says? & compare title, summary, and tags with the game \\
Is the content suitable for this Store's audience? & play past the title
  screen; read every string \\
Is the attribution honest? & authors, licence, acknowledgements; borrowed art
  and music \\
Is anyone's privacy at risk? & names, photographs, contact details in text or
  art \\
Is it a new version or a duplicate? & search the catalogue for the identifier
  and the title \\
\bottomrule
\end{tabularx}
\captionof{table}{An editorial checklist.}
\end{center}

A rejection should always come with a reason and a way forward, sent by
whatever channel the class uses; the Store records the decision, not the
conversation. Editors can also correct or withdraw what is already published:
the \kbtn{Cartridges} page lets them edit the title, summary, and tags of a
published version, or delete a variant or a whole version.

\begin{prgtry}
Run a peer-review session. Each student submits a cartridge to the class Store;
each is then made editor for \emph{someone else's} submission and must write
three sentences before deciding: what the game does, one thing that works well,
and one thing that must change before publication.
\end{prgtry}

\section{The administrator: running a Store}

\subsection{Deployment}

The Store is a small Python web service with an SQLite database and a data
directory. The supported way to run it is a container.

\lstinputlisting[style=bash,
  caption={First deployment of a classroom Store.}
]{scripts/store/admin_deploy.sh}

\begin{prgwarn}
A fresh Store has an administrator named \code{admin} whose password is
\code{password}. Change it before the first pupil connects. A Store also
refuses to start without a \code{SECRET\_KEY}; generate a long random one and
keep it with the deployment's secrets, because it signs every login session.
\end{prgwarn}

\subsection{Settings}

The \kbtn{Setup} page, visible to administrators only, holds everything that
gives a Store its identity and policy: its name, tagline, logo, colours, and
custom style; the default role of new users; publishing settings; whether and
how the Store advertises itself by mDNS as \code{\_prg32store.\_tcp.local.};
the SMTP server used to send registration links; and the optional external
login providers (OpenID Connect and SAML~2.0 from the page, LDAP from the
environment). Without SMTP, registration links are written to the server log,
which is acceptable only for a demonstration.

\subsection{People}

\begin{center}
\begin{tabularx}{\linewidth}{@{}lX@{}}
\toprule
Page & Purpose \\
\midrule
\kbtn{Users} & create, edit, and delete users; set a role; reset a password;
  assign groups \\
\kbtn{Groups} & create, rename, and delete groups; \code{editors} is the group
  that grants review rights \\
\code{/admin/roles} & the fixed roles and how many users hold each \\
\kbtn{Cartridges} & edit descriptive metadata or delete versions (shared with
  editors) \\
\kbtn{Backup} & download a complete backup; restore one \\
\bottomrule
\end{tabularx}
\captionof{table}{Administration pages \citep{prg32storeapi}.}
\end{center}

A sound arrangement for a course is: one or two administrators (the lecturer
and a technician); the teaching assistants in the \code{editors} group;
students as ordinary registered users. Do not make every editor an
administrator: the right to approve a game should not include the right to
read the user list or restore a backup.

\subsection{Backup, restore, and the network}

The \kbtn{Backup} page produces a single ZIP containing a snapshot of the
database and the data directory: settings, users, groups, cartridges, pending
submissions, scores, statistics, and metrics. The same page restores one.
Alternatively, stop the service and archive the data directory.

\lstinputlisting[style=bash,
  caption={A cold backup of the Store's data directory.}
]{scripts/store/admin_backup.sh}

\begin{center}
\begin{tabularx}{\linewidth}{@{}lX@{}}
\toprule
Check & Detail \\
\midrule
Reachability & the host firewall admits TCP 5080; boards and laptops share the
  network \\
Discovery & \code{/.well-known/prg32-store.json} returns addresses that are
  reachable \emph{from the boards} \\
Multiplayer & a reverse proxy, if any, forwards WebSocket traffic to
  \code{/api/multiplayer} \\
Transport & beyond a closed classroom network, serve the Store over HTTPS
  behind a reverse proxy \\
Tokens & rotate shared classroom tokens at the end of each term \\
Backups & scheduled, stored elsewhere, and restored once as a test \\
\bottomrule
\end{tabularx}
\captionof{table}{An operations checklist.}
\end{center}

\section{Publishing from the Construction Kit}

Chapter~\ref{ch:kids3} ended with a child pressing \kbtn{Publish}. What
happened is now easy to describe. The Kit generated C from the blocks, built
the two architecture variants, wrote a manifest, zipped a bundle, and sent it
to the Store with the token stored in the teacher's publish profile. The Store
created a pending submission in that teacher's name, and the teacher, wearing
the editor's hat, verified it. A seven-year-old's game travelled exactly the
road of Figure~\ref{fig:storeflow}.

\section{Publishing as a scholarly habit}

It is worth saying why a first-year course should bother with any of this. A
catalogue entry with an identifier, a version, named authors, a licence, and a
review record is, in miniature, a publication in the sense the rest of a degree
will use the word. Students who publish a cartridge have practised versioning,
attribution, licensing, peer review, and the difference between a private
artefact and a public one, on an object they care about.

\begin{prgcheck}
You have understood this chapter if you can: name the four roles and one thing
only each can do; explain the difference between a cartridge, a bundle, and a
game; say which manifest fields can never be changed after submission; create
and revoke a publication token; submit a bundle and find it in the pending
queue; and state the first two things an administrator must do on a new Store.
\end{prgcheck}
